\documentclass[11pt,a4paper]{article}

% 支持中文
\usepackage{xeCJK}
\usepackage[utf8]{inputenc}
\usepackage{geometry}
\usepackage{titlesec}
\usepackage{enumitem}
\usepackage{hyperref}

% 页面设置
\geometry{left=1.2cm, right=1.2cm, top=1.2cm, bottom=1.2cm}
\pagestyle{empty}

% 标题格式设置
\titleformat{\section}{\large\bfseries}{}{0em}{}[\titlerule]
\titlespacing*{\section}{0pt}{1.2ex}{1ex}

\begin{document}

% --- 个人信息 ---
\begin{center}
    {\LARGE \textbf{杨奥}} \\
    \vspace{0.8ex}
    满族 | 2006.02 | 共青团员 \\
    电话：139-2197-0469 | 邮箱：\href{mailto:1744393991@qq.com}{1744393991@qq.com} \\
    GitHub: \href{https://github.com/evil0knight}{github.com/evil0knight}
\end{center}

% --- 教育背景 ---
\section{教育背景}
\textbf{江苏大学} \hfill 2024.09 -- 至今 \\
物联网工程（本科） \hfill 
\begin{itemize}[leftmargin=1.5em, itemsep=0pt, topsep=2pt]
    \item \textbf{核心课程：} 408 计算机基础（数据结构/组成原理/操作系统/网络）、电路分析、数模电（考研难度标准）、Arduino 快速原型开发。
\end{itemize}

% --- 实习经历 ---
\section{实习经历}
\textbf{南京长城医疗器械有限公司} \hfill \textbf{嵌入式开发实习生} \\
\textit{项目一：MRI 环境下数码管显示驱动系统开发} \hfill 2024.12 -- 2025.01
\begin{itemize}[leftmargin=1.5em, itemsep=0pt]
    \item \textbf{背景挑战：} 针对核磁共振强磁环境下电容屏失效问题，负责设计高抗干扰的数码管替代方案。
    \item \textbf{技术实现：} 基于 \textbf{和泰 HT32} 芯片完成 \textbf{TM1640} 驱动 11 位数码管的软件 IIC 底层驱动；在 USART 接口集成 \textbf{TVS 二极管} 并完成 \textbf{RS232} 电平转换以满足医疗 EMC 检测标准。
    \item \textbf{工程实践：} 独立利用逻辑分析仪调试并修正了原理图引脚连点缺陷，通过微调 IIC 时序解决了复杂硬件环境下的通讯协议稳定性问题。
\end{itemize}

\vspace{1ex}

\textbf{宁波双马水控项目} \hfill \textbf{硬件/软件负责人} \\
\textit{项目二：433MHz 无线水控超低功耗接收系统} \hfill 2025.01 -- 2025.02
\begin{itemize}[leftmargin=1.5em, itemsep=0pt]
    \item \textbf{背景挑战：} 为最高 100V 电压的水流发电系统开发具备极端低功耗性能的无线控制模块。
    \item \textbf{功耗优化：} 深度调优 \textbf{FT61FC4} 芯片与 433 模块休眠逻辑，将系统平均电流从 \textbf{21mA 压缩至 45uA}，实现功耗 \textbf{467 倍} 的大幅下降。
    \item \textbf{系统设计：} 采用 60V 固态电容解决高压波动滤波；实现无线协议解包算法，支持物理层前导码唤醒与多机地址校验。
\end{itemize}

% --- 项目经历 ---
\section{项目经历}

\textbf{平衡轮腿 —— 全国大学生智能车竞赛} \hfill \textbf{项目负责人} \\
\begin{itemize}[leftmargin=1.5em, itemsep=0pt]
    \item \textbf{自驱攻坚：} 大一期间在\textbf{无代码传承、无学长老师指导}的极端背景下，依靠极强自学能力独立负责软件架构，实现“一人承担两人份”开发工作。
    \item \textbf{核心技术：} 独立编写基于 \textbf{大津法图像分割} 与 \textbf{种子搜线算法} 的识别模块；应用 \textbf{VMC 五连杆控制理论} 与 \textbf{三级串级 PID} 算法，攻克了轮腿机器人在复杂地形下的自平衡与跳跃难题。
    \item \textbf{项目成果：} 达成 \textbf{1.5m/s} 稳定车速，获全国大学生智能车竞赛华东赛区 \textbf{二等奖}，系\textbf{本校同届该赛道历史最好成绩}。
\end{itemize}

\vspace{4ex} % 进一步拉大间距，确保两个大项目视觉区分度极高

\textbf{超级电容管理系统 —— Robomaster 机甲大师赛} \hfill \textbf{硬件负责人} \\
\begin{itemize}[leftmargin=1.5em, itemsep=0pt]
    \item \textbf{技术突破：} 针对**团队战队 5 年未能攻克**的核心瓶颈，在大二转岗硬件后仅用 3 个月即完成两版硬件方案迭代，成功复刻并优化了业界主流方案。
    \item \textbf{硬件实现：} 独立设计基于 \textbf{H 桥自动升降压} 电路，实现高精度 \textbf{ADC 采样} 与 \textbf{CAN 总线} 通信，通过算法实现了功率的实时闭环控制（70W $\to$ 50W）。
    \item \textbf{产出影响：} \textbf{彻底填补了团队在超级电容电力管理领域的技术空白}，大幅提升了机器人运动爆发力。
\end{itemize}

% --- 专业技能 ---
\section{专业技能}
\begin{itemize}[leftmargin=1.5em, itemsep=0pt]
    \item \textbf{嵌入式开发：} 精通 C 语言，熟练使用 STM32 / HT32 / FT61FC4 / TC377；深入理解 PID、VMC 及基础图像处理算法。
    \item \textbf{硬件设计：} 熟练使用 \textbf{AD13 (Altium Designer)} 与嘉立创 EDA；掌握 Buck/Boost 电源设计、EMC 防护及器件选型。
    \item \textbf{工程能力：} 具备独立使用\textbf{逻辑分析仪}翻译通讯协议的能力；拥有从 0 到 1 的全栈开发与复杂硬件调试经验。
\end{itemize}

\end{document}